\title{Controlling Text-to-Image Diffusion by Orthogonal Finetuning}

\begin{abstract}
Large-scale text-to-image diffusion models have demonstrated remarkable proficiency in creating photorealistic imagery from textual inputs. A significant open question remains: how can we efficiently steer or condition these robust models to handle diverse downstream tasks? To this end, we present Orthogonal Finetuning~(OFT), a systematic approach for refining text-to-image diffusion models for specialized applications. Distinct from prior techniques, OFT is theoretically guaranteed to maintain hyperspherical energy, a property capturing pairwise neuron relations on the unit hypersphere. Our analysis indicates that maintaining this hyperspherical energy is vital to retaining the semantic generative capacity of text-to-image diffusion models. Additionally, to bolster the stability of the finetuning process, we introduce Constrained Orthogonal Finetuning (COFT), which supplements the original method with a radius constraint on the hypersphere. We particularly address two central finetuning scenarios: (i) subject-driven generation, which entails producing images of a specific subject in novel settings based on a small set of subject images alongside a text prompt, and (ii) controllable generation, wherein the model is extended to accommodate auxiliary control signals. Experimental results substantiate that our OFT methodology surpasses existing baselines in terms of generation fidelity and finetuning efficiency.
\end{abstract}

\section{Introduction}

Recent advancements in text-to-image diffusion models~\cite{saharia2022photorealistic,ramesh2022hierarchical,rombach2022high} have resulted in state-of-the-art performance for generating high-quality images under textual control. However, purely text-based conditioning remains inherently ambiguous and may fall short in delivering precise and detailed regulation of image synthesis. To address these limitations, our work concentrates on two pivotal categories of text-to-image generation tasks:

\begin{itemize}[leftmargin=*,nosep]

    \item \textbf{Subject-driven generation}~\cite{ruiz2023dreambooth}: Provided with a limited number of images featuring a particular subject, the challenge is to synthesize new depictions of the same subject situated in varying contexts, guided by a textual prompt.
    \item \textbf{Controllable generation}~\cite{zhang2023adding,mou2023t2i}: When furnished with an extra control input (\emph{e.g.}, canny edge maps, semantic segmentation masks), the objective is to produce images that are aligned both with the external control signal and the text instruction.
\end{itemize}

The central challenge in both tasks revolves around devising strategies to finetune text-to-image diffusion models while maintaining the generative strengths acquired during pretraining. The key requirements for an effective finetuning approach can be summarized as follows: (1) \emph{training efficiency}, defined by minimizing both the number of adjustable parameters and the overall training duration; and (2) \emph{generalizability preservation}, entailing the retention of the model’s capacity for high-quality and diverse generation. Standard finetuning procedures typically involve either incrementally updating existing neuron weights at a low learning rate (\emph{e.g.}, \cite{ruiz2023dreambooth}) or integrating lightweight modules with separate, re-parameterized weights (\emph{e.g.}, \cite{hulora2022,zhang2023adding}). While straightforward, neither method inherently ensures the preservation of generative performance obtained during pretraining. Additionally, there remains a gap in theoretical understanding for constructing optimal finetuning strategies or determining crucial hyperparameters such as the number of epochs required. A major obstacle stems from the absence of robust metrics to assess how well the pretrained generative abilities are retained post-finetuning. Contemporary approaches largely presume that reducing the Euclidean distance between the finetuned and pretrained models suffices to uphold pretrained capabilities. This presumption also underlies the prevalent choice of extremely low learning rates in finetuning. However, although this strategy may sometimes be effective, we contend that simply measuring Euclidean distance does not adequately reflect the extent of semantic retention. Thus, a more principled, structural metric for assessing the divergence between the finetuned and original models could substantially enhance both generative preservation and the reliability of the finetuning process.

Building on empirical findings that hyperspherical similarity serves as an effective encoder of semantic information~\cite{liu2017deep,liu2018decoupled,chen2020angular}, we leverage hyperspherical energy~\cite{liu2018learning} to quantify the inter-neuronal relationships within a layer. Specifically, hyperspherical energy is calculated as the aggregate of hyperspherical similarities (\emph{e.g.}, cosine similarities) between every pair of neurons in the same layer, offering a metric for neuron uniformity distributed over the unit hypersphere~\cite{liu2021learning}. We posit that effective finetuning should yield minimal change in hyperspherical energy relative to the pretrained state. A straightforward approach would be to introduce a regularization term to enforce consistent hyperspherical energy throughout finetuning; however, this strategy does not necessarily ensure an optimal minimization of the energy difference. As an alternative, we exploit a key invariance property of hyperspherical energy: applying a common orthogonal transformation across all neurons leaves the pairwise hyperspherical similarities invariant. Drawing from this invariance, we introduce Orthogonal Finetuning~(OFT), a methodology that adapts large text-to-image diffusion models for specific downstream tasks while preserving their hyperspherical energy. The principal concept involves learning a shared orthogonal transformation at the layer level, thereby maintaining the pairwise angular relationships among neurons. This procedure can also be conceptualized as a reparameterization of the neurons' coordinate system within the layer. Furthermore, acknowledging that maintaining a smaller Euclidean distance between the finetuned and pretrained models facilitates better retention of pretrained knowledge, we extend this approach to Constrained Orthogonal Finetuning~(COFT), a variant that enforces the updated neurons to remain on a hypersphere of fixed radius centered at their pretrained positions.

The rationale behind utilizing orthogonal transformations in neuron finetuning is partly informed by insights from the 2D Fourier transform, where an image can be represented in terms of its magnitude and phase spectra. The phase spectrum, which encodes the angular relationship between inputs and the basis, largely determines the semantic content. For instance, reconstructing an image by pairing its original phase spectrum with a randomized magnitude spectrum can retain the core semantics, underscoring the significance of angular information. This observation highlights that altering neuron orientations fundamentally changes the semantic attributes of generated images—a key objective in both subject-driven and controllable image synthesis. Nonetheless, if neuron directions are modified without constraint, the generative abilities acquired during pretraining can be compromised. To mitigate this, we propose maintaining the angles among neurons, motivated by the hypothesis that these angular relationships are vital to the knowledge representations within neural networks. Based on this perspective, we advocate for applying a layer-wise, shared orthogonal transformation that preserves the hyperspherical energy within each layer, thereby maintaining the underlying geometry.

Additionally, our approach is informed by the methodology of orthogonal over-parameterized training~\cite{liu2021orthogonal}, which involves training classification models from scratch through orthogonal transformations applied to randomly initialized networks. Random initializations result in low hyperspherical energy on average, and the main objective in~\cite{liu2021orthogonal} is to keep this energy minimal throughout training, which has been shown to enhance generalization for classification tasks~\cite{liu2018learning,lin2020regularizing}. It has been demonstrated that orthogonal transformations provide enough capacity to effectively train classifiers. In contrast, our focus is on adapting text-to-image diffusion models to improve control and enhance generation quality in downstream tasks. It is important to clarify the distinctions between OFT and~\cite{liu2021orthogonal} in two main respects. Firstly, whereas~\cite{liu2021orthogonal} targets minimizing hyperspherical energy, OFT aims to maintain it at the pretrained level to prevent disrupting the pretrained model’s inherent structure, as excessive minimization could damage the original semantic framework in diffusion models. Secondly, OFT introduces constraints to reduce divergence from the pretrained configuration, giving rise to its constrained formulation, whereas no such restrictions are enforced in~\cite{liu2021orthogonal}. Ultimately, effective finetuning requires a balance between flexibility and stability; our OFT method is designed to strike this balance. Our contributions can be summarized as follows:

\begin{itemize}[leftmargin=*,nosep]

    \item We introduce a new finetuning approach termed Orthogonal Finetuning, designed to enhance the controllability of text-to-image diffusion models. Additionally, to bolster the method’s robustness, we present a constrained version that restricts the angular change relative to the original pretrained model.
    \item In contrast to traditional finetuning techniques, OFT updates model parameters in a manner that provably maintains hyperspherical energy—a property we empirically identify as crucial for retaining the generative semantics from the pretrained model.
    \item The OFT framework is evaluated across two application domains: subject-driven generation and controllable generation. Through a broad experimental analysis, we demonstrate that OFT yields marked improvements over earlier methods in generation fidelity, training efficiency, and stability throughout finetuning. Notably, OFT attains strong results even when trained on substantially fewer images and epochs, underscoring its sample efficiency.
    \item For controllable generation, we devise a novel control consistency metric to assess controllability. This metric operates by predicting the control signal from the generated output and comparing it against the original control input, thus providing further evidence of OFT’s efficacy in delivering strong controllability.
\end{itemize}

\section{Related Work}

\textbf{Text-to-image diffusion models}. Substantial advancements~\cite{saharia2022photorealistic,ramesh2022hierarchical,rombach2022high,gu2022vector,nichol2022glide} have been realized in generating images from text, primarily propelled by innovations in diffusion-based generative modeling~\cite{ho2020denoising,song2021score,song2021denoising,dhariwal2021diffusion} as well as vision-language learning~\cite{su2020vlbert,lu2019vilbert,radford2021learning,wang2021simvlm,li2022blip,li2023blip,alayrac2022flamingo,schuhmann2022laion}. Both GLIDE~\cite{nichol2022glide} and Imagen~\cite{saharia2022photorealistic} focus on modeling in pixel space; GLIDE jointly trains the text encoder alongside a diffusion prior using paired data, whereas Imagen utilizes a frozen, pretrained text encoder. In contrast, Stable Diffusion~\cite{rombach2022high} and DALL-E2~\cite{ramesh2022hierarchical} model diffusion processes within latent space. Stable Diffusion leverages VQ-GAN~\cite{esser2021taming} for constructing a visual codebook latent space, while DALL-E2 employs CLIP~\cite{radford2021learning} to create a shared embedding space for image and text representations. Beyond diffusion models, approaches based on generative adversarial networks~\cite{reed2016generative,zhang2017stackgan,xu2018attngan,li2019controllable} and autoregressive architectures~\cite{ramesh2021zero,ding2021cogview,wu2022nuwa,yuscaling} have also contributed to the field of text-to-image generation. The OFT technique is inherently compatible with different diffusion models for finetuning and is not model-specific.

\textbf{Subject-driven generation}. In order to minimize changes to the central subject, \cite{avrahami2022blended,nichol2022glide} use user-provided masks as auxiliary conditioning inputs. Inversion approaches~\cite{choi2021ilvr,dhariwal2021diffusion,ramesh2022hierarchical,gal2022image} allow for context alteration while preserving subject identity. Furthermore, \cite{hertz2022prompt} facilitates both localized and global edits in the absence of explicit input masks. These prior solutions, however, struggle to retain detailed identity attributes. Pivotal Tuning~\cite{roich2022pivotal} finetunes a generator centered around an inverted latent code, incorporating an extra regularization mechanism for identity preservation. In a similar vein, \cite{nitzan2022mystyle} trains a personalized generative face prior using a collection of facial images from an individual. Meanwhile, \cite{casanova2021instance} is able to produce diverse instance variations, though at the expense of potentially missing out on finer instance-specific traits. DreamBooth~\cite{ruiz2023dreambooth} adapts the text-to-image diffusion model to new subjects by learning with a limited number of subject images, a custom token, and losses for reconstruction and class prior preservation. Our method, OFT, utilizes the DreamBooth setup but replaces straightforward parameter updates with orthogonal transformations, rather than relying solely

\textbf{Controllable generation}. Image-to-image translation falls under the broader category of controllable generation, with earlier approaches largely relying on conditional generative adversarial networks~\cite{isola2017image,zhu2017toward,wang2018high,park2019semantic,choi2018stargan}. More recent developments have seen the adoption of diffusion models in this domain~\cite{saharia2022palette,voynov2022sketch,wang2022pretraining}. Notably, ControlNet~\cite{zhang2023adding} introduces a mechanism for steering pretrained diffusion models by finetuning them with supplemental control signals, attaining impressive controllability outcomes. T2I-Adapter~\cite{mou2023t2i}, a parallel effort, similarly enhances control in generated imagery via targeted finetuning of pretrained diffusion models. In line with the frameworks outlined in \cite{zhang2023adding,mou2023t2i}, we employ OFT to adapt pretrained diffusion models, demonstrating superior controllability while requiring less training data and fewer finetuned parameters. Crucially, OFT does not entail additional computational cost at inference.

\textbf{Model finetuning}. The practice of adapting large, pretrained models to downstream applications through finetuning is increasingly prevalent~\cite{devlin2018bert,brown2020language,he2022masked}. Among the various adaptation strategies (\emph{e.g.}, \cite{hulora2022,houlsby2019parameter,pfeiffer2020adapterfusion}) widely explored in natural language processing, LoRA~\cite{hulora2022} stands out as particularly relevant to OFT, as it introduces low-rank parameter updates during finetuning. In contrast, OFT adopts a layer-wise, shared orthogonal transformation, adjusting neuron weights multiplicatively. This approach theoretically preserves inter-neuronal angular relationships within each layer, leading to improved model stability.

\section{Orthogonal Finetuning}

\subsection{Why Does Orthogonal Transformation Make Sense?}

We begin by examining the motivation for utilizing orthogonal transformations during finetuning of text-to-image diffusion models. Our analysis is organized around two fundamental questions: (1) what is the rationale for adjusting neuron angles (i.e., their directional properties) during finetuning, and (2) what advantages does orthogonal transformation offer for such angular adjustments?

To address the first question, we take guidance from prior empirical findings in \cite{liu2018decoupled,chen2020angular}, which report that the angular disparity between features is an effective indicator of semantic divergence. SphereNet~\cite{liu2017deep} demonstrates that constraining all neuron activations to reside on a unit hypersphere during network training not only preserves model capacity but can also improve generalization, suggesting that the orientation of neuron vectors sufficiently encodes salient information from input data. To further elucidate the role of neuron angles, we present a controlled experiment in Figure~\ref{fig:toy_ae}, in which a conventional convolutional autoencoder is trained on a dataset of flower images. During training, the standard inner product is employed to compute feature maps, where $z$ represents the output of the convolutional filter $\bm{w}$ acting on the input patch $\bm{x}$. At inference, we assess three approaches for constructing the feature map: (a) the inner product, consistent with training; (b) the extraction of only magnitude components; and (c) the utilization of purely angular components. As depicted in Figure~\ref{fig:toy_ae}, reconstructing the input images is nearly perfect when relying solely on angular information, whereas the magnitude alone proves uninformative. It is important to note that the cosine-based (angular) activation (c) is not used during training; the model is trained strictly with the inner product. This result underscores that it is principally the neuron directions—rather than magnitudes—that encode the semantic essence of the data. Thus, adjusting these directions may offer a particularly effective lever for semantic manipulation of images.

Concerning the second question, a straightforward strategy to update neuron orientations involves collectively rotating or reflecting all neurons within the same layer, inherently invoking orthogonal transformations. While more adaptive angular manipulations (allowing distinct rotations per neuron) might offer increased flexibility, our observations suggest that orthogonal transformations optimally balance adaptability and structural coherence. Additionally, \cite{liu2021orthogonal} has demonstrated the adequacy of orthogonal transformations in the context of learning in neural networks. To substantiate this position, we conduct a controlled generation experiment within the ControlNet~\cite{zhang2023adding} framework, as illustrated in Figure~\ref{fig:ortho}. Our findings reveal that the standard OFT consistently delivers stable performance and precise control after sufficient training epochs (epoch 20). Conversely, the variant of OFT lacking the orthogonality condition fails to synthesize plausible images or exhibit control. These results reinforce the necessity of the orthogonality constraint within OFT.

\subsection{General Framework}

Traditional finetuning methods often rely on gradient descent with a low learning rate to modify model parameters (or specific layers), which encourages the updated weights to remain close to those of the original pretrained model. This process can be interpreted as implicitly minimizing $\thickmuskip=2mu \medmuskip=2mu \| \bm{M} - \bm{M}^0\|$, where $\bm{M}$ represents the current weights after finetuning and $\bm{M}^0$ stands for the initial pretrained weights. While this constraint intuitively ensures stability, it may still grant too much flexibility for adaptation in large-scale models. To mitigate this, LoRA enforces an explicit low-rank condition on the update such that $\thickmuskip=2mu \medmuskip=2mu \text{rank}(\bm{M} - \bm{M}^0)=r'$, with $r'$ being a small predefined value. Unlike LoRA, OFT applies a restriction based on the neuron-wise pairwise similarity through hyperspherical energy, formulated as $\thickmuskip=2mu \medmuskip=2mu \| \text{HE}(\bm{M})-\text{HE}(\bm{M}^0) \|=0$, where $\text{HE}(\cdot)$ is the hyperspherical energy function over the weight matrix. 

For illustration, consider a fully connected layer $\thickmuskip=2mu \medmuskip=2mu \bm{W}=\{\bm{w}_1,\cdots,\bm{w}_n\}\in\mathbb{R}^{d\times n}$ where each column $\bm{w}_i\in\mathbb{R}^{d}$ denotes the $i$-th neuron's weights and $\bm{W}^0$ refers to its initialization. The output $\thickmuskip=2mu \medmuskip=2mu \bm{z}\in\mathbb{R}^n$ is given by $\thickmuskip=2mu \medmuskip=2mu \bm{z}=\bm{W}^\top\bm{x}$, with $\thickmuskip=2mu \medmuskip=2mu \bm{x}\in\mathbb{R}^d$ as the input. OFT can thus be interpreted as minimizing the difference in hyperspherical energy between finetuned and original weights:

\begin{equation}\label{eq:oft_principle}
\footnotesize
\min_{\bm{W}} \left\| \text{HE}(\bm{W})-\text{HE}(\bm{W}^0)\right\|~~\Leftrightarrow~~\min_{\bm{W}} \bigg{\|} \sum_{i\neq j} \|\hat{\bm{w}}_i-\hat{\bm{w}}_j\|^{-1}-\sum_{i\neq j} \|\hat{\bm{w}}^0_i-\hat{\bm{w}}^0_j\|^{-1}\bigg{\|}
\end{equation}

where $\thickmuskip=2mu \medmuskip=2mu \hat{\bm{w}}_i:=\bm{w}_i/\|\bm{w}_i\|$ denotes the normalized vector for neuron $i$, and hyperspherical energy is defined as $\thickmuskip=2mu \medmuskip=2mu \text{HE}(\bm{W}):= \sum_{i\neq j} \|\hat{\bm{w}}_i-\hat{\bm{w}}_j\|^{-1}$ for the layer $\bm{W}$. It is straightforward to note that the lowest value possible for Eq.~\eqref{eq:oft_principle} is zero, which occurs if and only if $\bm{W}$ is obtained from $\bm{W}^0$ by an orthogonal transformation—i.e., $\thickmuskip=2mu \medmuskip=2mu \bm{W}=\bm{R}\bm{W}^0$ for some orthogonal $\thickmuskip=2mu \medmuskip=2mu \bm{R}\in\mathbb{R}^{d\times d}$ (with determinant $1$ or $-1$ for rotation or reflection, respectively). The core principle of OFT, then, is to finetune the network by learning an orthogonal transformation per layer, shared by all neurons in that layer. In this framework, OFT seeks

\begin{equation}\label{eq:oft_general}
    \footnotesize
    \bm{z} = \bm{W}^\top\bm{x} = (\bm{R}\cdot\bm{W}^0)^\top\bm{x},~~~\text{s.t.}~\bm{R}^\top\bm{R} = \bm{R}\bm{R}^\top = \bm{I}
\end{equation}

Here, $\bm{W}$ refers to the weight matrix obtained after OFT finetuning, while $\bm{I}$ stands for the identity matrix. An overview of the OFT method can be seen in Figure~\ref{fig:block}. To parallel the zero-initialization principle used in LoRA, OFT must begin finetuning from the pretrained parameter state $\bm{W}^0$. This requirement is met by initializing the orthogonal matrix $\bm{R}$ as an identity matrix, thereby preserving the pretrained weights at the start of finetuning. Ensuring the orthogonality of $\bm{R}$ can be accomplished using differentiable orthogonalization techniques such as those proposed by \cite{liu2021orthogonal,lezcano2019cheap}. We explore computational approaches that efficiently maintain this orthogonality constraint.

\subsection{Efficient Orthogonal Parameterization}\label{sect:eff_ortho}

Conventional orthogonalization procedures like the Gram-Schmidt process, despite being differentiable, are generally impractical for large-scale use due to high computational costs~\cite{liu2021orthogonal}. For improved computational efficiency, we utilize Cayley parameterization to obtain an orthogonal matrix. Specifically, we define the orthogonal matrix as $\thickmuskip=2mu \medmuskip=2mu \bm{R}=(\bm{I}+\bm{Q})(I-\bm{Q})^{-1}$, where $\bm{Q}$ is a skew-symmetric matrix such that $\thickmuskip=2mu \medmuskip=2mu \bm{Q}=-\bm{Q}^\top$. While this approach boosts efficiency, it restricts the constructed matrices to the special orthogonal group (having determinant $1$); in our experiments, this constraint does not result in practical drawbacks. Even using Cayley parameterization, representing $\bm{R}$ for high-dimensional $d$ may still be inefficient in terms of parameters. To alleviate this, we represent $\bm{R}$ as a block-diagonal matrix composed of $r$ blocks, yielding the structure:

\begin{equation}
    \footnotesize
    \bm{R} = \text{diag}(\bm{R}_1,\bm{R}_2,\cdots,\bm{R}_r) = \begin{bmatrix}
    \bm{R}_1\in \text{O}(\frac{d}{r}) &  &\\
    &\ddots & \\
    & & \bm{R}_r\in \text{O}(\frac{d}{r})
    \end{bmatrix}\in \text{O}(d)
\end{equation}

Here, $\text{O}(d)$ represents the group of $d$-dimensional orthogonal matrices, with $\thickmuskip=2mu \medmuskip=2mu \bm{R}\in\mathbb{R}^{d\times d}$ and $\thickmuskip=2mu \medmuskip=2mu \bm{R}_i\in\mathbb{R}^{d/r\times d/r}$ for each $i$ being orthogonal. In the special case where $\thickmuskip=2mu \medmuskip=2mu r=1$, the block-diagonal structure degenerates into a standard fully unconstrained orthogonal matrix. The parameter count for a $d \times d$ orthogonal matrix is $\thickmuskip=2mu \medmuskip=2mu d(d-1)/2$, implying an overall parameter complexity of $\mathcal{O}(d^2)$. For an orthogonal matrix partitioned into $r$ blocks, the parameters reduce to $\thickmuskip=2mu \medmuskip=2mu d(d/r-1)/2$, leading to a complexity of $\thickmuskip=2mu \medmuskip=2mu \mathcal{O}(d^2/r)$. Additionally, parameter sharing among block matrices (that is, imposing $\thickmuskip=2mu \medmuskip=2mu\bm{R}_i=\bm{R}_j$ for all $i \neq j$) can further compress the parameter space, lowering the complexity to $\mathcal{O}(d^2/r^2)$. Importantly, regardless of such efficiency measures, the composite matrix $\bm{R}$ remains orthogonal, ensuring hyperspherical energy preservation is unaffected.

To compare OFT with LoRA in parameter efficiency, note that LoRA uses $r'$ as a low-rank parameter and requires $\thickmuskip=2mu \medmuskip=2mu r'(d+n)$ trainable parameters. If both $r$ and $r'$ scale linearly with $d$ (for example, $\thickmuskip=2mu \medmuskip=2mu r=r'=\alpha d$ with some $\thickmuskip=2mu \medmuskip=2mu 0<\alpha\leq 1$), LoRA’s parameter count grows as $\mathcal{O}(d^2+dn)$, whereas OFT scales as $\mathcal{O}(d)$. To make this comparison concrete, consider a $128\times 128$ weight matrix: LoRA would use $2,048$ parameters at $\thickmuskip=2mu \medmuskip=2mu r'=8$, while OFT uses $960$ parameters for $\thickmuskip=2mu \medmuskip=2mu r=8$, provided block sharing is not utilized.

\subsection{Constrained Orthogonal Finetuning}

The OFT framework can be made more restrictive by limiting how far the adapted model can deviate from the original pretrained model. Constrained Orthogonal Finetuning (COFT) achieves this by modifying the forward computation as follows:

\begin{equation}\label{eq:coft_general}
    \footnotesize
    \bm{z}=\bm{W}^\top\bm{x}=(\bm{R}\cdot\bm{W}^0)^\top\bm{x},~~~\text{s.t.}~\bm{R}^\top\bm{R}=\bm{R}\bm{R}^\top=\bm{I},~\norm{\bm{R}-\bm{I}}\leq \epsilon
\end{equation}

Here, the constraints enforce both orthogonality and proximity (with $\epsilon$ as the bound) of $\bm{R}$ to the identity matrix. The orthogonality property can be maintained by employing the Cayley parameterization, as described in Section~\ref{sect:eff_ortho}. However, embedding the $\epsilon$-proximity condition into the Cayley-parameterized framework is a challenging problem. To further interpret the Cayley transform, we utilize the Neumann series to approximate $\thickmuskip=2mu \medmuskip=2mu \bm{R}=(\bm{I}+\bm{Q})(I-\bm{Q})^{-1}$ by $\thickmuskip=2mu \medmuskip=2mu \bm{R}\approx\bm{I}+2\bm{Q}+\mathcal{O}(\bm{

Because the series converges under the operator norm, we can equivalently impose the restriction $\thickmuskip=2mu \medmuskip=2mu\|\bm{R}-\bm{I}\|\leq\epsilon$ within the Cayley transform. This translates to the constraint $\thickmuskip=2mu \medmuskip=2mu \|\bm{Q}-\bm{0}\|\leq\epsilon'$, where $\bm{0}$ denotes a zero matrix and $\epsilon'$ serves as a separate hyperparameter (distinct from $\epsilon$). Enforcement of this new constraint on $\bm{Q}$ is straightforward using projected gradient descent. To guarantee that the orthogonal matrix $\bm{R}$ is initialized as the identity, $\bm{Q}$ is set to the zero matrix at initialization. COFT therefore integrates two clear constraints: minimizing hyperspherical energy deviation, and explicitly bounding deviations from the pretrained parameters. The latter is typically only implicit in other finetuning methods, whereas COFT formulates it directly. Although OFT performs well, COFT demonstrates even greater finetuning stability, owing to its explicit deviation control. Figure~\ref{fig:eps_coft} illustrates how variations in $\epsilon$ influence COFT’s performance: larger $\epsilon$ permits more adaptability, making the COFT-finetuned model similar to one finetuned with OFT; as $\epsilon$ decreases, the resulting COFT model increasingly retains the behavior of the original pretrained text-to-image diffusion model.

\subsection{Re-scaled Orthogonal Finetuning}

To augment OFT, we introduce a variant in which each neuron’s activation is scaled by a learnable magnitude parameter. This stems from the observation that such rescaling does not impact the hyperspherical energy, as magnitudes are ultimately normalized. Formally, the revised forward computation is given by $\thickmuskip=2mu \medmuskip=2mu\bm{z}=(\bm{R}\bm{W}^0\bm{D})^\top\bm{x}$\footnote{\textbf{Errata}: The NeurIPS camera ready version erroneously states the re-scaled OFT forward pass as $\thickmuskip=2mu \medmuskip=2mu\bm{z}=(\bm{D}\bm{R}\bm{W}^0)^\top\bm{x}$. The implementation is correct and results in Appendix~\ref{app:rescaled} are unaffected.} where $\thickmuskip=2mu \medmuskip=2mu \bm{D}=\text{diag}(s_1,\cdots,s_n)\in\mathbb{R}^{n\times n}$, a diagonal matrix parameterized by positive values $s_1,\cdots,s_n$. Unlike standard OFT (Eq.~\eqref{eq:oft_general}) where only $\bm{R}$ is learned, both the scaling matrix $\bm{D}$ and orthogonal matrix $\bm{R}$ are trained here. The introduction of $\bm{D}$ modestly increases parameter count, yet offers enhanced expressive power over plain OFT. Although we retain standard OFT in main experiments to isolate the effect of the orthogonal transformation, empirical findings (see Appendix~\ref{app:rescaled}) indicate re-scaled OFT often yields superior performance.

\section{Intriguing Insights and Discussions}

\textbf{OFT is architecture-independent}. In principle, OFT can be adapted to any neural network architecture. For Transformers, the prevailing use of LoRA involves modifying attention weights~\cite{hulora2022}, so for consistency in evaluation, we limit OFT’s application to attention layers during experiments. OFT’s design is not limited to fully connected layers: its block-diagonal $\bm{R}$ structure affords interpretable benefits in convolutional contexts, differentiating it from LoRA. By assigning one block per input channel, each block targets a specific neuron channel, akin to learning depthwise convolution kernels~\cite{chollet2017xception}. If all blocks in $\bm{R}$ are identical and shared, then neurons across channels are subjected to the same orthogonal transformation. We present exploratory results on applying OFT to convolutional layers in Appendix~\ref{app:convolution}.

\textbf{Connection to LoRA}. LoRA achieves stability during adaptation by incorporating a low-rank matrix, thereby limiting substantial alterations to the pretrained weights. Conversely, OFT regulates modifications to the pretrained weight matrix through an orthogonal (full-rank) transformation, ensuring preservation of information learned during pretraining. The OFT forward computation can be reformulated as $\thickmuskip=2mu \medmuskip=2mu \bm{z}=(\bm{R}\bm{W}^0)^\top\bm{x}=(\bm{W}^0+(\bm{R}-\bm{I})\bm{W}^0)^\top\bm{x}$, in which $\thickmuskip=2mu \medmuskip=2mu (\bm{R}-\bm{I})\bm{W}^0$ functions similarly to the low-rank updates in LoRA. When $\bm{W}^0$ maintains full rank and $\thickmuskip=2mu \medmuskip=2mu \bm{R}-\bm{I}$ is of low rank, OFT can also be interpreted as performing a low-rank weight adjustment. Both approaches include mechanisms for parameter reduction: LoRA utilizes a tunable rank $r'$, whereas OFT employs a block-diagonal parameter $r$ to decrease the count of parameters subject to training. Notably, LoRA and OFT represent two fundamentally different strategies for parameter-efficient tuning: LoRA relies on low-rank representations, whereas OFT leverages structural sparsity in the form of block-diagonal orthogonal transformations.

\textbf{Why OFT converges faster}. First, as evidenced in Figure~\ref{fig:toy_ae}, altering neuron directions proves to be the most potent means to shift network semantics—precisely the aspect that OFT is architected to manipulate. Additionally, OFT’s parameterization constrains updates to remain on a smooth hypersphere manifold, effectively enabling finetuning in a more favorable optimization landscape. Empirical results in \cite{liu2021orthogonal} further corroborate these theoretical insights.

\textbf{Why not minimize hyperspherical energy}. Distinct from the approach in \cite{liu2021orthogonal}, our objective does not involve minimizing hyperspherical energy. In classification contexts, optimal performance is attained when neurons encode unique, non-redundant features. Minimizing hyperspherical energy enforces uniform distribution of neuron directions over the hypersphere, a constraint that may disrupt the informative priors encoded in pretrained models, thus rendering it unsuitable for finetuning tasks.

\textbf{Balancing flexibility and regularity during finetuning}. Our analysis reveals a fundamental trade-off between the adaptability and stability of finetuning approaches. While standard finetuning methods offer the greatest adaptability, they often compromise stability and can result in model collapse. OFT, despite its conceptual simplicity, achieves an effective compromise by maintaining the pairwise angles between neuron vectors, thus mediating between flexibility and regularization. The block-diagonal approach can be interpreted as applying even stricter regularization to the orthogonal transformation.

\textbf{Inference efficiency is maintained}. In contrast to ControlNet, the OFT framework imposes no extra computational load during inference on the finetuned model. During inference, the learned orthogonal matrix $\bm{R}$ is directly multiplied with the pretrained weights $\bm{W}^0$, producing a resultant weight matrix $\thickmuskip=2mu \medmuskip=2mu \bm{W}=\bm{R}\bm{W}^0$. As a result, the inference performance remains on par with the original pretrained model.

\section{Experiments and Results}

\textbf{Experimental configurations}. Our empirical evaluation is based on the Stable Diffusion v1.5~\cite{rombach2022high} text-to-image generation model. To ensure fairness, we randomly select images generated from each tested method. For subject-driven tasks, our protocol follows DreamBooth~\cite{ruiz2023dreambooth}. For experiments involving control signals, the evaluation setup generally follows ControlNet~\cite{zhang2023adding} and T2I-Adapter~\cite{mou2023t2i}. For fair comparison with LoRA, OFT or COFT is restricted to the same layers used by LoRA. Further methodological details and expanded results are available in Appendix~\ref{app:settings}.

\subsection{Subject-driven Generation}

\textbf{Experimental setup}. DreamBooth~\cite{ruiz2023dreambooth} and LoRA~\cite{hulora2022} are used as baseline approaches. All competing methods are trained using the same loss function as specified in DreamBooth. The training procedure for DreamBooth and LoRA follows their respective original protocols, with hyperparameters set according to best practices suggested in their papers. Additional findings can be found in Appendix~\ref{app:settings},~\ref{app:coft_oft},~\ref{app:more_qualitative},~\ref{app:failure}.

\textbf{Finetuning stability and convergence}. We begin by assessing both the stability and convergence rate of finetuning for DreamBooth, LoRA, OFT, and COFT, with the corresponding findings depicted in Figure~\ref{fig:teaser} and Figure~\ref{fig:db_convergence}. Our analysis reveals that COFT and OFT enable stable finetuning for the diffusion model throughout training. At the 400-iteration mark, DreamBooth and OFT variants have achieved robust control over the generation process, whereas LoRA exhibits difficulty maintaining subject identity. Upon reaching 2000 iterations, DreamBooth suffers from mode collapse, leading to suboptimal image outputs, and LoRA fails to reproduce the yellow shirt, substituting yellow fur instead. Conversely, OFT and COFT retain consistent and reliable controllability of the generated results. These outcomes underscore the capacity of the OFT framework to converge swiftly without sacrificing stability, particularly in the context of subject-driven tasks. Importantly, the models’ robustness to the specific choice of iteration number mitigates the need for meticulous adjustment per subject. Both OFT and COFT allow practitioners to select a large iteration count by default. Additionally, with an appropriate choice of $\epsilon$ in COFT, setting both the learning rate and number of iterations becomes largely straightforward.

\textbf{Quantitative comparison}. Using the evaluation methodology described in \cite{ruiz2023dreambooth}, we quantitatively assess subject fidelity (using DINO~\cite{caron2021emerging} and CLIP-I~\cite{radford2021learning}), text prompt alignment (CLIP-T~\cite{radford2021learning}), as well as sample diversity (LPIPS~\cite{zhang2018unreasonable}). Specifically, CLIP-I computes the mean pairwise cosine similarity between CLIP embeddings of generated images and corresponding real images. DINO performs a similar operation utilizing ViT S/16 DINO embeddings instead. CLIP-T quantifies the average cosine similarity between CLIP embeddings of text prompts and generated outputs. Furthermore, we calculate the mean LPIPS cosine similarity for samples generated from identical prompts and subjects. The quantitative findings, reported in Table~\ref{table:dreambooth_final}, demonstrate that both COFT and OFT exceed the performance of DreamBooth and LoRA on the DINO and CLIP-I measures by a substantial margin and deliver competitive or superior outcomes in prompt fidelity and diversity assessments. To ensure reliability, each algorithm is repeatedly finetuned using the same pretrained model with 30 different random seeds. This analysis confirms that the OFT framework not only delivers enhanced convergence and stability, but also achieves consistently superior final results.

\textbf{Qualitative comparison}. To provide clearer insight into the advantages offered by OFT, we present several randomly selected samples of subject-driven image generation in Figure~\ref{fig:dreambooth_final}. For consistency and fairness, each example is produced using the identical finetuned model for all methods, ensuring that no prompt is individually tailored to improve any specific outcome. For each approach, the version of the model with the highest validation CLIP score is utilized. Upon examining the outcomes in Figure~\ref{fig:dreambooth_final}, it is evident that both OFT and COFT excel at maintaining semantic integrity of the subject, whereas LoRA frequently fails to retain subject characteristics (\emph{e.g.}, in the case of the bowl, LoRA does not preserve the subject’s identity at all). Additionally, OFT and COFT allow for significantly more precise text-driven manipulation of the generation process, unlike DreamBooth, which, while often managing to preserve subject identity, does not always produce images that align well with the intended text prompts (\emph{e.g.}, the first instance in the bowl row). Overall, the qualitative analysis highlights that OFT surpasses competitors by delivering superior subject retention while providing enhanced controllability through text. Furthermore, OFT’s performance remains stable regardless of the number of iterations, whereas both DreamBooth and LoRA display notable declines under such circumstances. Further qualitative illustrations are available in Appendix~\ref{app:more_qualitative}. Additionally, Appendix~\ref{app:human} presents a human assessment that corroborates OFT’s superiority.

\subsection{Controllable Generation}

\textbf{Settings}. We adopt ControlNet~\cite{zhang2023adding}, T2I-Adapter~\cite{mou2023t2i}, and LoRA~\cite{hulora2022} as comparative baselines. The evaluation focuses on three difficult controllable generation benchmarks described in the main text: translating Canny edges to images~(C2I) using the COCO dataset~\cite{lin2014microsoft}, converting segmentation maps to images~(S2I) based on the ADE20K dataset~\cite{zhou2017scene}, and mapping landmarks to facial images~(L2F) using CelebA-HQ~\cite{karras2018progressive,xia2021tedigan}. Every method is employed to finetune Stable Diffusion~(SD) v1.5 for 20 training epochs on each dataset. Expanded results and further discussion can be found in Appendix~\ref{app:more_qualitative}, \ref{app:more_control_task}, \ref{app:failure}.

\textbf{Convergence}. To assess how quickly each method converges on the L2F task, we analyze ControlNet, T2I-Adapter, LoRA, and COFT, employing both numerical and visual evaluations. For the quantitative aspect, we measure the mean $\ell_2$ distance between the set of control face landmarks and those predicted by the models. As presented in Figure~\ref{fig:face_convergence}, we display the facial landmark errors across varying epochs for models trained with different approaches. It is clear from these results that both COFT and OFT exhibit much more rapid convergence compared to the baselines. Notably, LoRA requires 20 epochs to match the performance level that our OFT framework attains by only the 8-th epoch. Importantly, OFT and COFT maintain a comparable number of trainable parameters to LoRA, which is substantially less than what ControlNet uses, yet they reach optimal performance much faster than previous methods. Further evidence for the swift convergence of OFT appears in Figure~\ref{fig:teaser}, where the example on the right demonstrates OFT's superior data efficiency: OFT achieves convergence using just 5\% of the ADE20K dataset, outpacing ControlNet and LoRA. For the qualitative evaluation, our comparison centers on OFT, COFT, and ControlNet, as ControlNet is the method whose landmark errors most closely resemble ours. As illustrated in Figure~\ref{fig:face_convergence_qual}, OFT and COFT not only achieve stable convergence, but also ensure that the synthesized face pose increasingly aligns with the control landmarks throughout training. By contrast, ControlNet displays a delayed emergence of controllability—only after the 8-th epoch—which echoes the abrupt convergence described in \cite{zhang2023adding}. The consistent and gradual convergence behavior observed in OFT significantly simplifies practical deployment, as its smoother and more predictable learning dynamics make it easier to identify effective hyperparameters for the OFT framework.

\textbf{Quantitative comparison}. To assess controllability in generation tasks, we propose a control consistency metric, which involves extracting the control signal from a generated image and matching it to the input control signal. For the C2I task, performance is measured using the IoU and F1 scores, while for the S2I task, we calculate mean IoU as well as both mean and overall accuracy. For the L2F task, evaluation relies on the mean $\ell_2$ distance between predicted and actual control landmarks. Additional information about the evaluation metrics can be found in Appendix~\ref{app:settings}. Each competing method is tested using optimally selected hyperparameters. As reported in Table~\ref{table:control_gen}, both OFT and COFT demonstrate superior accuracy and robustness in control compared to alternative methods. Notably, adapter-based techniques (\emph{e.g.}, T2I-Adapter and ControlNet) exhibit slower convergence and yield suboptimal outcomes. LoRA achieves improved results over ControlNet on S2I, but underperforms on C2I and L2F tasks. Overall, LoRA’s performance remains capped at a modest level, even after extensive hyperparameter optimization. In contrast, our OFT framework continues to improve as we increase the number of trainable parameters, indicating the ceiling has yet to be reached. It is worth highlighting that the introduced quantitative protocol for measuring control in generation is a key contribution, enabling rigorous evaluation of the finetuned models—constrained primarily by the inherent accuracy of the reference segmentation or detection models.

\textbf{Qualitative comparison}. We perform a qualitative assessment of OFT and COFT against prevailing state-of-the-art techniques such as ControlNet, T2I-Adapter, and LoRA. As illustrated by the randomly generated samples in Figure~\ref{fig:control_final}, OFT and COFT are distinguished by their capacity to produce images of high realism and fidelity, while simultaneously providing precise control over the generation process. In the S2I task, it is evident that LoRA is unable to adhere to the constraints provided by the segmentation maps, whereas ControlNet, OFT, and COFT consistently generate images that respect the input structure. Compared to ControlNet, both OFT and COFT generate images with sharper details and more plausible geometric configurations, despite employing significantly fewer parameters. For the C2I task, OFT and COFT excel in synthesizing semantically consistent outputs from rudimentary Canny edge inputs, a feat where T2I-Adapter and LoRA fall short. In the L2F task, our method demonstrates superior pose control over generated facial images, successfully handling difficult facial orientations. Across all three control scenarios, OFT and COFT deliver images of higher qualitative value than current state-of-the-art baselines, attesting to the robust controllability enabled by our OFT approach. For a more thorough qualitative evaluation, we include additional examples of all three control scenarios in Appendix~\ref{app:control_exp}. Further, Appendix~\ref{app:more_control_task} illustrates the versatility of OFT in various other control settings, such as dense pose to human body, sketch-to-image, and depth-to-image tasks.

\section{Concluding Remarks and Open Problems}

Inspired by the pivotal role that angular relationships among neurons play in defining visual semantics, we introduce a straightforward yet robust fine-tuning approach—orthogonal finetuning—to enhance control over text-to-image diffusion models. This strategy specifically addresses two tasks: subject-driven generation and controllable generation. In comparison with prior techniques, OFT achieves superior controllability and more stable finetuning outcomes, all while utilizing fewer parameters. Notably, OFT does not increase inference time, ensuring that the resulting models are practical for real-world deployment.

Nevertheless, OFT gives rise to several notable open questions. Firstly, OFT ensures orthogonality by employing Cayley parametrization, which relies on computing a matrix inverse—this can restrict the scalability of the approach. While we mitigate this issue by adopting block diagonal parametrization, developing methods to accelerate this matrix inversion in a differentiable manner remains unresolved. Secondly, OFT holds distinctive promise in compositionality: orthogonal matrices generated through multiple OFT fine-tuning tasks can be multiplied to yield another orthogonal matrix. It remains an open question whether these matrices consistently retain knowledge across all downstream tasks. Lastly, the extent to which OFT achieves parameter efficiency is tied to the use of block diagonal structures, which can introduce bias and constrain flexibility. Identifying more effective, less biased means to increase parameter efficiency stands out as an important avenue for future research.

\end{document}